\documentclass{article}
\usepackage{amsmath, amssymb, amsthm}
\usepackage[margin=1in]{geometry}
\usepackage{hyperref}

\newtheorem{theorem}{Theorem}
\newtheorem{lemma}[theorem]{Lemma}
\newtheorem{proposition}[theorem]{Proposition}
\newtheorem{corollary}[theorem]{Corollary}
\newtheorem{conjecture}[theorem]{Conjecture}
\theoremstyle{definition}
\newtheorem{definition}[theorem]{Definition}
\newtheorem{remark}[theorem]{Remark}
\newtheorem*{observation}{Observation}

\title{BMS Hybrid-Basis Chain and the Staircase $\Pi_q^{(n)}$:\\
A Documented Misalignment, with a Sharpened Gap for $\sigma_2$-G1 at $n=6$}
\author{Clio}
\date{2026-04-30}

\begin{document}

\maketitle

\begin{abstract}
We document an attempted closure of $\sigma_2$-G1 at $n=6$ via the
Brauner--Meszaros--Saliola (BMS) hybrid-basis chain
(arXiv:2602.19508). The PROVE-session strategy assumed a structural
analogy between (i) the $n=4$ BMS verification of Apr 30 and (ii) the
staircase Bott--Samelson element $\Pi_q^{(n)}$, conjecturing that
$\Pi_q^{(6)}$ in the Kazhdan--Lusztig basis is supported on the
Bruhat interval $[e, s_1 s_2 s_3 s_4 s_5]$. We show by direct
computation that this analogy is \emph{false} already at $n=4$: the
Apr 30 BMS verification was for the length-3 element
$(1+T_1)(1+T_2)(1+T_3)$, whereas the staircase
$\Pi_q^{(4)}$ has length~$\binom{4}{2} = 6$ and full KL support
across all 24 elements of $S_4$. The BMS chain does not, by itself,
bound the KL support of $\Pi_q^{(n)}$. We then explain why the
remaining BMS technology --- factorization of the KL change-of-basis
through parabolic-restriction blocks --- still does not deliver
$\sigma_2$-G1: the BMS factorization is a change of basis, not an
endomorphism factorization, and $\sigma_2$ is not multiplicative
along basis-change steps. We close by sharpening the trace-square
gap of \texttt{2026-04-28-sigma2-G1-n6-trace-square.tex}: the open
question is the palindromic positivity of the difference
$(\sigma_1^2 - \operatorname{tr}(\Pi^2))/(2q^{a+1}(1+q)^{2b})$, for
which BMS supplies no further structure.
\end{abstract}

\section{Setup and the PROVE-session conjecture}

Let $H_q = H_q(S_n)$ with $T_i^2 = (q-1)T_i + q$ and $q = v^2$. For
$n \geq 2$ define the \emph{staircase reduced word}
\[
\mathbf{i}_n \;=\; (1; \; 2,1; \; 3,2,1; \; \ldots; \; n-1, n-2, \ldots, 1),
\qquad \ell(\mathbf{i}_n) = \binom{n}{2} =: N.
\]
This word is reduced for $w_0 \in S_n$. Set
\[
\Pi_q^{(n)} \;=\; \prod_{j=1}^{N} \bigl(1 + T_{i_j}\bigr) \;\in\; H_q,
\]
the Bott--Samelson element along $\mathbf{i}_n$. Let $V_\lambda$ be
Hoefsmit's $q$-seminormal Specht module for $\lambda \vdash n$, and
write $\sigma_k(\lambda) = e_k(\text{eigenvalues of }\Pi_q^{(n)}|_{V_\lambda})$.

The \emph{$\sigma_2$-G1 property} for $V_\lambda$ is: $\sigma_2(\lambda)
= q^{a_{\lambda,2}}(1+q)^{b_{\lambda,2}} P_{\lambda, 2}(q)$ with
$P_{\lambda, 2} \in \mathbb{Z}_{\geq 0}[q]$ palindromic.

\medskip
The PROVE-session plan (file \texttt{/home/clio/state/PROVE.md},
``$\sigma_2$-G1 at $n=6$ even shapes via BMS hybrid-basis chain'')
proposes closure via the BMS chain
$\emptyset \subset \{s_1\} \subset \{s_1, s_2\} \subset \{s_1, s_2,
s_3\} \subset \{s_1, \ldots, s_4\} \subset S_6$. Step 2 of that plan
reads, verbatim:
\begin{quote}
\emph{Step 2. Compute $\Pi_q$ in $\operatorname{KL}(S_6)$ basis.
Conjecturally support is on Bruhat interval $[e, s_1 s_2 s_3 s_4 s_5]$
(the increasing word) by analogy with $n=4$ --- verify computationally.}
\end{quote}

\section{The misalignment}

The ``analogy with $n=4$'' refers to the Apr 30 in-container
computation in
\texttt{/home/clio/projects/scratch/2026-04-30-bms-hybrid-chain-n4/},
whose \texttt{RESULT.md} reports:
\begin{quote}
\emph{$\Pi_q$ in $\operatorname{KL}(S_4)$ basis: only 8 of 24 coefficients
non-zero; pure powers of $(1-v)$; supported on Bruhat interval
$[e, s_1 s_2 s_3]$.}
\end{quote}

A direct read of \texttt{bms\_chain.py} in that folder shows that the
``$\Pi_q$'' computed there is
\[
\Pi^{\text{BMS}} \;:=\; (1 + T_1)(1 + T_2)(1 + T_3),
\qquad \ell = 3,
\]
\emph{not} the staircase $\Pi_q^{(4)}$, whose length is $\binom{4}{2}
= 6$:
\[
\Pi_q^{(4)} \;=\; (1+T_1)(1+T_2)(1+T_1)(1+T_3)(1+T_2)(1+T_1).
\]

\begin{observation}[Computational, n=4]
\label{obs:n4-full-support}
Expand $\Pi_q^{(4)}$ in the Kazhdan--Lusztig basis $\{C_w\}_{w \in S_4}$.
All $24$ coefficients are non-zero; the support is therefore the entire
Bruhat order of $S_4$, not the interval $[e, s_1 s_2 s_3]$.
\end{observation}

This was verified by the script
\texttt{/home/clio/projects/scratch/2026-04-30-bms-n6-pi-kl/staircase\_pi\_in\_KL.py},
which reports for $n=4$:
\begin{verbatim}
staircase word for n=4: [1, 2, 1, 3, 2, 1], length 6
Pi_q^(4) has 24 nonzero T-basis terms (out of 24)
Pi_q^(4) has 24 nonzero KL-basis terms
\end{verbatim}
A representative non-zero coefficient at $w = w_0$:
$[\Pi_q^{(4)}]_{C_{w_0}} = 1$. A representative non-zero coefficient
at $w = e$:
$[\Pi_q^{(4)}]_{C_e} = -(v-1)\,(v^6 - 5v^5 + 9v^4 - 8v^3 + 6v^2 - 2v + 1)$.

\medskip
The same script for $n=3$ gives staircase
$\Pi_q^{(3)} = (1+T_1)(1+T_2)(1+T_1)$ with all $6$ KL coefficients
non-zero. The simpler element $\Pi^{\text{BMS}} = (1+T_1)(1+T_2)$ in
$H_q(S_3)$ has only the four coefficients on $[e, s_1 s_2]$ non-zero
--- which is the analogue of the $8/24$ figure at $n=4$ for $\Pi^{\text{BMS}}$.

\begin{corollary}\label{cor:no-bruhat-bound}
The Bruhat-interval support claim of \texttt{PROVE.md} Step~2 cannot
be inferred from the Apr~30 $n=4$ verification. The two computations
concern different elements of the Hecke algebra: the Apr~30
computation diagonalizes $\Pi^{\text{BMS}} = \prod_{i=1}^{n-1}(1+T_i)$,
\emph{not} the staircase $\Pi_q^{(n)}$.
\end{corollary}

\section{Why BMS does not deliver $\sigma_2$-G1 anyway}

Even setting aside the support claim, we now explain why the BMS
chain does not close $\sigma_2$-G1 by the proof strategy of
\texttt{PROVE.md} Step~3.

\subsection{BMS factorizes a change of basis, not an endomorphism}

Fix the chain
$\emptyset = J_0 \subset J_1 \subset \cdots \subset J_r = S$ of subsets
of simple reflections. BMS defines a hybrid basis $\mathrm{TC}^{J_i}$
of $H_q(W)$ for each $i$ and shows that the change of basis
$M_i: \mathrm{TC}^{J_{i-1}} \to \mathrm{TC}^{J_i}$ is block-diagonal:
$M_i = \mathrm{id}_{W^{J_{i-1}}} \otimes K_{W_{J_i}/W_{J_{i-1}}}$,
where $K_{W_{J}/W_{J'}}$ is a parabolic-Hecke KL change-of-basis matrix.
The full KL change-of-basis is the product
\[
B \;=\; M_r \cdots M_2 \cdot M_1 : \mathrm{T} \longrightarrow \mathrm{C}.
\]

The matrix $\Pi_q^{(n)}$ acting on $V_\lambda$ does \emph{not} factor
along this chain. The $M_i$ are basis-change matrices on $H_q$
itself, not on $V_\lambda$, and even if one passes to $V_\lambda$ via
the Specht-module quotient, the BMS factorization does not lift to a
factorization
\[
\Pi_q^{(n)}|_{V_\lambda} \;\stackrel{?}{=}\; A_r \cdots A_2 \cdot A_1
\]
by endomorphisms $A_i \in \mathrm{End}(V_\lambda)$. Indeed
$\Pi_q^{(n)}|_{V_\lambda}$ has a fixed Jordan structure determined by
$\lambda$ that ignores the BMS chain entirely.

\subsection{$\sigma_2$ is not multiplicative}

For matrices $A, B \in \mathrm{End}(V_\lambda)$,
\[
\sigma_2(A B) \;\neq\; \sigma_2(A)\,\sigma_2(B)
\quad \text{in general.}
\]
Concretely, for $\dim V_\lambda = m$ with eigenvalues $\mu_1, \ldots,
\mu_m$ of $A$ and $\nu_1, \ldots, \nu_m$ of $B$, the eigenvalues of
$AB$ are not the products $\mu_i \nu_i$ unless $A, B$ commute and
share an eigenbasis. The PROVE.md ``Open technical hinge'' anticipates
this:
\begin{quote}
\emph{$\sigma_2$ is not multiplicative on matrix product (it's the
second elementary symmetric polynomial of eigenvalues). The right
sub-multiplicativity is on the spectrum of the chain's output, not
the chain's matrix product itself.}
\end{quote}
The PROVE.md proposal then suggests recovering $\sigma_2$-G1 via
``positivity of the $M_i$ matrix entries plus $\sigma_2$-G1 of the
fiber rep at smaller $n$.'' But this conflates two different
operators: the $M_i$ act on $H_q(W)$ as a basis change, while
$\sigma_2$-G1 is a property of $\Pi_q^{(n-1)}$ acting on $V_\mu$ for
$\mu \vdash n-1$. There is no direct mechanism converting positivity
of $M_i$ into positivity of $\sigma_2(\Pi_q^{(n)}|_{V_\lambda})$.

\subsection{The branching primitive is the survivor; BMS does not refine it}

What \emph{does} restrict $\Pi_q^{(n)}$ to the parabolic
$H_q(W_{J_i})$ in a useful way is the fact that $\Pi_q^{(n)}$
contains $\Pi_q^{(n-1)}$ as a left factor:
\[
\Pi_q^{(n)} \;=\; \Pi_q^{(n-1)} \cdot \prod_{i=1}^{n-1}(1 + T_{n-1-i+1} \cdots T_1)^{?}
\]
--- no, this is wrong: the staircase decomposes as
\[
\Pi_q^{(n)} \;=\; \Pi_q^{(n-1)} \cdot (1+T_{n-1})(1+T_{n-2})\cdots (1+T_1),
\]
i.e.\ the last $n-1$ factors form the ``new row'' of the staircase.
This decomposition \emph{is} a useful endomorphism factorization, and
it underlies the parity-asymmetry analysis of
\texttt{2026-04-25-parity-asymmetry.md}: at odd $n$ the parabolic
subgroup $H_n = \langle s_1, s_3, \ldots, s_{n-1}\rangle$ sits inside
$S_{n-1}$, giving a clean branching reduction; at even $n$ this fails.
The BMS chain re-expresses the parabolic structure of the KL basis,
but does not refine the branching primitive: it does not provide a
new endomorphism factorization of $\Pi_q^{(n)}$, and so it cannot
close the parity-asymmetry gap.

\section{Sharpened gap for $\sigma_2$-G1 at $n=6$}

\subsection{The trace-square reformulation}

By Newton's identity,
\begin{equation}\label{eq:newton}
2 \sigma_2(\lambda) \;=\; \sigma_1(\lambda)^2 \;-\; \operatorname{tr}\bigl(\Pi_q^{(n)\, 2}|_{V_\lambda}\bigr).
\end{equation}
By $\sigma_1$-G1 (Apr 24, file \texttt{2026-04-24-sigma1-G1.tex}):
$\sigma_1(\lambda) = \sum_{k} n_k(\lambda)(1+q)^k q^{(N-k)/2}$ with
$n_k \geq 0$. Squaring,
\begin{equation}\label{eq:sigma1sq}
\sigma_1(\lambda)^2 \;=\; \sum_{l} \widetilde N_l(\lambda)(1+q)^l q^{(2N - l)/2},
\quad \widetilde N_l \geq 0.
\end{equation}
The element $\Pi_q^{(n)\, 2}$ is itself a Bott--Samelson element along
the doubled word $\mathbf{i}_n \cdot \mathbf{i}_n$ (length $2N$), to
which $\sigma_1$-G1 again applies, giving
\begin{equation}\label{eq:trsq}
\operatorname{tr}\bigl(\Pi_q^{(n)\, 2}|_{V_\lambda}\bigr)
\;=\; \sum_{l} M_l(\lambda)(1+q)^l q^{(2N-l)/2},
\quad M_l \geq 0.
\end{equation}
Subtracting and dividing by 2:
\begin{equation}\label{eq:sigma2}
\sigma_2(\lambda) \;=\; \tfrac{1}{2} \sum_{l} \bigl(\widetilde N_l(\lambda) - M_l(\lambda)\bigr)
(1+q)^l q^{(2N-l)/2}.
\end{equation}

\subsection{What BMS would have had to deliver}

For $\sigma_2$-G1 to follow from the trace-square reformulation, one
would need either
\begin{enumerate}
\item[(C1)] $\widetilde N_l(\lambda) \geq M_l(\lambda)$ for all $l$
(the strong G1-cone bound), or
\item[(C2)] a structural cancellation grouping the terms in
\eqref{eq:sigma2} into manifestly G1 blocks.
\end{enumerate}
\texttt{2026-04-29-sigma2-G1-KL-cone-decomposition.tex} verifies (C1)
for $\lambda = (5,1)$ and $(4,2)$ at $n=6$ and refutes it for
$\lambda = (3,2,1)$ (negative coefficient $-3$ at $l=4$). The
BMS-chain hope was that (C2) would emerge from the parabolic
factorization of the KL change-of-basis. Section~3 above shows that
no such cancellation can be extracted from the BMS factorization
alone, since BMS factorizes a basis change, not the endomorphism
$\Pi_q^{(n)}|_{V_\lambda}$.

\subsection{The remaining gap}

The closure of $\sigma_2$-G1 at $n=6$ for all rank-$\geq 2$ shapes
$(5,1), (4,2), (3,2,1)$ is established by \emph{two} mechanisms:
KL-cone (Apr 29) for $(5,1)$ and $(4,2)$; atom decomposition (Apr 29,
\texttt{2026-04-29-sigma2-G1-321-atom-decomp.tex}) for $(3,2,1)$.
What remains open is a \emph{single uniform mechanism} that handles
all three. Sections 2 and 3 above rule out the BMS chain as such a
mechanism.

\begin{conjecture}[Sharpened uniform-closure gap]\label{conj:gap}
There exists $\epsilon \in \mathbb{Z}_{\geq 0}$ such that the
difference
\[
\Delta(\lambda; q) \;:=\; \sigma_1(\lambda)^2 - \operatorname{tr}\bigl(\Pi_q^{(n)\, 2}|_{V_\lambda}\bigr)
\]
admits a factorization
\[
\Delta(\lambda; q) \;=\; 2 q^{a_\lambda + 1}(1+q)^{b_\lambda}\,
Q_\lambda(q)
\]
with $Q_\lambda \in \mathbb{Z}_{\geq 0}[q]$ palindromic, derived from
a structural invariant of $\Pi_q^{(n)}$ that is \emph{not} a parabolic
KL change-of-basis. Candidates: (i) the BS bimodule
$\mathrm{BS}(\mathbf{i}_n)$ in the Soergel category and its
Hochschild cohomology; (ii) the difference
$\mathrm{BS}(\mathbf{i}_n \cdot \mathbf{i}_n) - \mathrm{BS}(\mathbf{i}_n) \otimes \mathrm{BS}(\mathbf{i}_n)$ as a graded $H_q$-bimodule.
\end{conjecture}

\section{Computational verification record}

\subsection{Staircase $\Pi_q^{(4)}$ has full KL support}

Script: \texttt{2026-04-30-bms-n6-pi-kl/staircase\_pi\_in\_KL.py}. Output:
$\Pi_q^{(4)} = \sum_{w \in S_4} c_w(v) C_w$ with all $c_w \neq 0$. In
particular, $c_e = -(v-1)(v^6 - 5v^5 + 9v^4 - 8v^3 + 6v^2 - 2v + 1)$
and $c_{w_0} = 1$. Compare to $\Pi^{\text{BMS}} = (1+T_1)(1+T_2)(1+T_3)$
which has $c_w = 0$ for the $16$ elements $w \notin [e, s_1 s_2 s_3]$.

\subsection{$\sigma_1$-G1 base for the doubled word}

The doubled staircase word $\mathbf{i}_n \cdot \mathbf{i}_n$ is again
a reduced expression for some element of $S_n$? No: it has length
$2N$ which exceeds $\ell(w_0) = N$, so $\mathbf{i}_n \cdot \mathbf{i}_n$
is \emph{not} reduced, but the Bott--Samelson product
$\Pi_q^{(n)} \cdot \Pi_q^{(n)} = \Pi_q^{(n)\, 2}$ is still a positive
sum of $T_w$'s, and its trace decomposes by the same
$\sigma_1$-G1 mechanism (Theorem 1 of
\texttt{2026-04-24-sigma1-G1.tex}) applied to the BS element along
the non-reduced word $\mathbf{i}_n \cdot \mathbf{i}_n$.
This is verified at $n = 6$ in
\texttt{2026-04-28-sigma2-G1-n6-trace-square.tex}.

\section{Status update for the program}

\begin{itemize}
\item \textbf{$\sigma_2$-G1 at $n=6$:} closed shape-by-shape (KL-cone
for $(5,1), (4,2)$; atom decomposition for $(3,2,1)$). Uniform
closure remains open.
\item \textbf{BMS chain:} \emph{ruled out} as uniform-closure
mechanism, by Sections 2 and 3.
\item \textbf{Trace-square gap:} sharpened to Conjecture~\ref{conj:gap}.
The right structural object is plausibly the Soergel bimodule
$\mathrm{BS}(\mathbf{i}_n)$ and its Hochschild cohomology
(\texttt{2026-04-30-hh-bs-n4-conjecture.md}), \emph{not} the BMS hybrid
basis.
\end{itemize}

\section*{Files referenced}
\begin{itemize}
\item \texttt{/home/clio/state/PROVE.md} (this session's directive)
\item \texttt{/home/clio/projects/scratch/2026-04-30-bms-hybrid-chain-n4/RESULT.md}
\item \texttt{/home/clio/projects/scratch/2026-04-30-bms-hybrid-chain-n4/bms\_chain.py}
\item \texttt{/home/clio/projects/scratch/2026-04-30-bms-n6-pi-kl/staircase\_pi\_in\_KL.py}
\item \texttt{/home/clio/projects/proofs/2026-04-24-sigma1-G1.tex}
\item \texttt{/home/clio/projects/proofs/2026-04-28-sigma2-G1-n6-trace-square.tex}
\item \texttt{/home/clio/projects/proofs/2026-04-29-sigma2-G1-KL-cone-decomposition.tex}
\item \texttt{/home/clio/projects/proofs/2026-04-29-sigma2-G1-321-atom-decomp.tex}
\end{itemize}

\end{document}
